\documentclass[12pt]{article}

% \linespread{1.35} % for line distances, not change this

\usepackage[colorlinks=true, urlcolor=blue, linkcolor=blue]{hyperref}
\usepackage{amsmath, amssymb, amsfonts, amsthm}
\usepackage{breqn} % math related
\usepackage{paralist} % to put items in line
\usepackage[most]{tcolorbox}
\usepackage{fontawesome5}
\usepackage{xcolor}
\tcbuselibrary{skins,breakable,theorems}
\usepackage{xepersian}

\definecolor{QuestionBlue}{RGB}{30,90,170}
\definecolor{QuestionBack}{RGB}{245,248,255}
\definecolor{MyGreen}{RGB}{40,120,60}
\definecolor{MyGreenBack}{RGB}{243,250,243}
\emergencystretch=1em  % Add to preamble or before problematic section1

\newtcbtheorem[
    auto counter,
    number within=section
]{question}{سؤال}{
    enhanced,
    breakable,
    colback=QuestionBack,
    colframe=QuestionBlue,
    coltitle=white,
    fonttitle=\bfseries,
    title filled,
    attach boxed title to top right={
        xshift=-2mm,
        yshift=-2mm
    },
    boxed title style={
        colback=QuestionBlue,
        colframe=QuestionBlue,
        arc=2mm
    },
    arc=2mm,
    boxrule=0.8pt,
    left=2mm,
    right=2mm,
    top=2mm,
    bottom=2mm,
    before skip=12pt,
    after skip=12pt,
}{q}

\newtcolorbox{greenbox}[1]{
    enhanced,
    breakable,
    colback=MyGreenBack,
    colframe=MyGreen,
    coltitle=white,
    fonttitle=\bfseries,
    title={#1},
    title filled,
    attach boxed title to top right={
        xshift=-4mm,
        yshift=-2mm
    },
    boxed title style={
        colback=MyGreen,
        colframe=MyGreen,
        arc=2mm
    },
    arc=2mm,
    boxrule=0.8pt,
    left=2mm,
    right=2mm,
    top=2mm,
    bottom=2mm,
    before skip=10pt,
    after skip=10pt
}

\newenvironment{solution}
{
    \par\medskip
    \hrule
    \smallskip
    \noindent\textbf{راه‌حل.}
    \par\smallskip
}
{
    \par\smallskip
    \smallskip
    \hrule
    \medskip
}



% \settextfont{XB Roya}
\settextfont{Sharif 1.2 Regular}

\begin{document}

\begin{titlepage}

\centering

\vspace*{2cm}

{\Large\textbf{دانشگاه صنعتی شریف}}

\vspace{0.5cm}

{\large دانشکده مهندسی کامپیوتر}

\vspace{2cm}

{\Huge\bfseries گزارش پروژه دوم}

\vspace{0.5cm}

{\Huge رای‌گیری تاییدی چندمرحله‌ای}

\vspace{2.5cm}

{\large
\begin{tabular}{rc}
    \textbf{درس:} & نظریه بازی‌ها\\[0.3cm]
    \textbf{استاد:} & دکتر سمیرا حسین‌قربان\\[0.3cm]
    \textbf{دانشجو:} & ایلیا یزدانی
\end{tabular}
}

\vfill

{\large تابستان ۱۴۰۵}

\end{titlepage}

\tableofcontents

\newpage

\begin{abstract}
    در این گزارش به پرسش‌های مطرح شده در پروژه داده شده به طور کامل پاسخ می‌دهیم.
\end{abstract}

\section{رای‌دهنده}

\subsection*{پرسش‌های بخش اطلاعات کامل}

\begin{enumerate}
    \item {

        بله این پروفایل همواره یک تعادل نش است.
        
        در این سیستم امتیاز هر کاندید صرفا به رای های هر نفر بستگی دارد و هر چه در 
        مجموعه‌های تاییدی بیشتری قرار داشته باشد، 
        شانس او برای قرار گرفتن در کمیته نهایی افزایش
        پیدا می‌کند. در واقع برخلاف قاعده‌های
         دیگری که معرفی شده، کاندیدهای دیگر در کمیته نهایی در
        امتیاز یک کاندید تاثیری ندارند. با توجه به این موضوع، به طور شهودی
        مشخص است که یکی از بهترین استراتژی هر رای‌دهنده این است که تنها تمام کاندیدهای
        مورد تاییدش را در لیستش قرار دهد.

        برای اثبات دقیق این موضوع، به طور تدریجی بهترین پاسخ رای‌دهنده را به مجموعه
        کاندیدهای مورد تاییدش تبدیل می‌کنیم، به طوری که سود او کاهش نیابد. فرض کنید
        رای‌دهنده مورد بررسی 
        (برای راحتی او را ایلیا می‌نامیم)
        به مجموعه $B'$
        رای داده باشد در حالی که مجموعه‌ی کاندیدهای مورد تایید او $B$ 
        باشد.

        تا زمانی که کاندید $c'$ 
        در $B'$ 
        وجود دارد که در $B$ 
        نیست، آن را از $B'$ 
        خارج می‌کنیم. با اینکار تنها نماینده‌ای که امتیاز از دست می‌دهد 
        همان $c'$
        است. بنابراین تنها تغییر ممکن در کمیته نهایی خارج شدن 
        $c'$
        و ورود کاندید دیگری است که قطعا مطلوبیت ایلیا را کاهش نمی‌دهد.
        با انجام این روند خواهیم داشت
        $B' \subseteq B $
        .
        
        فرض کنید $c \in B$
        باشد که در $B'$
        نیست. با اضافه کردن این کاندید در $B'$
        بدون تغییر امتیاز باقی کاندیدها، تنها امتیاز $c$
        افزایش می‌یابد. بنابراین تنها تغییر ممکن در کمیته انتخاب شده، 
        این است که $c$ وارد شود و کاندیدی دیگر خارج شود. وارد شدن
        $c$
        مطلوبیت ایلیا را یک واحد افزایش می‌دهد و خارج شدن یک کاندید مطلوبیت ایلیا
        را حداکثر یک واحد کاهش می‌دهد. پس در این تغییر هم مطلوبیت ایلیا کم
        نمی‌شود. با انجام این روند روی تمام $c \in B$
        خواهیم داشت $B = B'$.

        پس ما با یک سری تغییر باعث شدیم بدون کم شدن از مطلوبیت ایلیا
        رای او صادقانه شود. در واقع:
        \[
        u_i(F(B, B_{-i})) \ge u_i(F(B',B_{-i}))
        \]

        بنابراین برای هر رای‌دهنده اثبات کردیم که بهترین پاسخ او به هر پروفایلی
        رای صادقانه است. این موضوع نتیجه می‌دهد که این پروفایل یک تعادل نش
        است.
        }
        \item {
            در این بازی، مانند تمام بازی‌ها، بهترین پاسخ رای‌دهنده، کنشی
            است که مطلوبیت او را بیشینه کند. 
            (با فرض اطلاع از استراتژی باقی رای‌دهنده‌ها)
            در واقع در این بازی بهترین پاسخ، بهترین مجموعه‌ای است که 
            یک نفر می‌تواند انتخاب کند که بیشترین کاندید موردتاییدش
            در کمیته نهایی قرار بگیرد.

            در سوال اول همین بخش ثابت کردیم که رای صادقانه، همیشه بهترین 
            پاسخ است. (در واقع در اثبات آن از اینکه استراتژی باقی 
            رای‌دهنده‌ها چیست استفاده نکردیم)
            پس جواب این سوال 
            "بله"
            است.
        }
        \item {
            همانطور که در پاسخ سوال اول همین بخش گفته شد، مهمترین نکته‌ای
            باعث جلوگیری از انحراف سودمند در این قاعده می‌شود، این
            است که هر کاندید به طور مستقل از باقی کاندیدها امتیازدهی
            می‌شود و انتخاب یک کاندید روی انتخاب کاندید دیگر اثری نمی‌گذارد.

            یکی دیگر از ویژگی‌های موثر در این موضوع این است تغییر
            رای یک بازیکن نتیجه را چندان تغییر نمی‌دهد. همانطور که بررسی کردیم
            هر یک جابجایی در رای یک نفر، نهایتا موجب یک تغییر در کمیته نهایی
            می‌شود.

            دو ویژگی دیگر ذکر شده ارتباط چندانی بر جلوگیری از انحراف ندارند.
            نحوه شکستن حالت تساوی، اگر معقول باشد باعث انحراف یا عدم انحراف نمی‌شود.
            فاصله زیاد میان کاندیدهای منتخب و غیرمنتخب هم تاثیر چندانی ندارد.

            به طور کلی چیزی که می‌تواند رای‌هنده را به رای غیرصادقانه سوق دهد،
            این است که با قرار ندادن یک یا چند کاندید موردتاییدش
            در مجموعه رای خود، باعث قرار گرفتن کاندیدهای بیشتری در کمیته نهایی شود.
            موضوعی که با این قاعده قابل انجام نیست.
        }
        \item {
            با یک مثال نشان می‌دهیم که در قاعده \lr{PAV}
            لزوما پروفایل صادقانه، تعادل نش نیست.
            فرض کنید 
            $5$ بازیکن،
            $4$ کاندید
            و $k = 3$
            باشد. همچنین مجموعه‌های موردتایید به صورت زیر باشد:
            \begin{align*}
                A_1 &= \{c_1, c_2 \} \\
                A_2 &= \{c_1, c_2, c_4\} \\
                A_3 &= \{c_1, c_3, c_4 \} \\
                A_4 &= \{c_1, c_3, c_4 \} \\
                A_5 &= \{c_1, c_3, c_4 \} \\
            \end{align*}
            در این بازی، پروفایل صادقانه و قاعده
            \lr{PAV}
            منجر به کمیته نهایی 
            $c_1, c_4, c_3$
            می‌شود.
            (امتیاز نهایی این کمیته $5 + \frac{4}{2} + \frac{3}{3} = 8$ است
            که بیشترین مقدار ممکن است)

            اما اگر بازیکن شماره $1$
            به جای رای صادقانه، تنها به $c_2$
            رای دهد، آنگاه خواهیم داشت:
            \begin{align*}
                Score_{PAV}(\{c_1, c_4, c_2\}) &= 4 + \frac{4}{2} + \frac{4}{3} = 7.33 \\
                Score_{PAV}(\{c_1, c_4, c_3\}) &= 4 + \frac{4}{2} + \frac{3}{3} = 7\\
            \end{align*}
            بنابراین کمیته نهایی
            $c_1, c_4, c_2$
            خواهد بود، که یعنی سود بازیکن اول برابر $2$
            خواهد شد. بنابراین این بازیکن انگیزه انحراف دارد و پروفایل صادقانه 
            تعادل نش نیست.

            همین مثال برای 
            \lr{Approval CC}
            هم نشان می‌دهد که در این قاعده نیز پروفایل صادقانه لزوما تعادل نش نیست.

            همانطور که توضیح داده شد، تفاوت بنیادین این قواعد با قاعده 
            اول این است که امتیاز برای هر کمیته تعریف می‌شود و انتخاب کاندیدها
            به همدیگر ربط دارد. مثلا در این بازی، بازیکن اول از آنجایی که می‌داند
            با توجه به رای دیگران کاندید $c_1$
            انتخاب می‌شود، می‌تواند با تک‌رای به $c_2$
            ارزش حضور این کاندید در کمیته نهایی را افزایش دهد و کاری کند او
            به جای 
            $c_3$ 
            قرار بگیرد. دلیل این موضوع این است که این دو قاعده تلاش می‌کنند کمیته نهایی
            تنها از رای اکثریت تشکیل نشود و گروه‌های کوچکتر که نسبت خوبی در جمعیت دارند
            هم نماینده داشته باشند.
        }
        \item {
            به طور کلی اثر این مقدار در استراتژی بازیکنان به پارامترهای زیادی وابسته
            است و نمی‌توان آن را به طور مستقل دقیق تحلیل کرد.
            اما می‌توان گفت برای مقادیر کوچک $k$
            انگیزه انحراف کمتر است. مثلا برای $k=1$
            پروفایل صادقانه قطعا یک تعادل نش است.

            هرچقدر 
            $k$
            بزرگتر باشد، احتمال انحراف بیشتر می‌شود. دلیل این موضوع این است که ممکن است
            برخی رای‌دهندگان از رای دادن به کاندیدهایی که مورد قبول اکثریت هستند
            پرهیز کنند و به‌جای آن تنها به کاندیدهای دیگر مورد تاییدشان رای دهند.
            (همانطور که در مثال سوال چهار این بخش دیدیم)
            مثلا برای
            $k=m-1$
            اگر رای‌دهنده‌ای 
            تنها حامی یک کاندید باشد می‌تواند با تک‌رای به این کاندید
            احتمال حضور او را در کمیته بیشتر کند و مطلوبیت رای‌دهنده افزایش پیدا کند.
        }
\end{enumerate}

\subsection*{پرسش‌های بخش اطلاعات ناقص}

\begin{enumerate}
    \item {
        برای یک رای‌دهنده، فضای نوع 
        $\Theta_i$
        مچموعه تمام زیرمجموعه‌های ممکن از مجموعه کاندیدهاست. مثلا اگر دو
        کاندید داشته باشیم:
        \[
        \Theta_i = \{\{c_1,c_2\}, \{c_1\}, \{c_2\}, \{\}\}
        \]
        چیزی که هر رای‌دهنده می‌داند مجموعه کاندیدهای مورد تایید خودش است.
        باور احتمالاتی هر رای‌دهنده هم برای 
        هر رای‌دهنده دیگر و مجموعه‌ای از کاندیدها مشخص می‌کند که 
        چقدر احتمال دارد مجموعه موردتایید آن رای‌دهنده، این مجموعه کاندید باشد.
        اما یک رای‌دهنده با اطمینان نمی‌داند که این مجموعه برای هر رای‌دهنده دیگر
        دقیقا چیست.

        مثلا در همان مثال 
        \(2\)
        کاندید، یک رای‌دهنده می‌داند کدام کاندیدها موردتایید
        خودش است، اما نمی‌داند رای‌دهنده دیگر کدام کاندیدها را قبول دارند.
    }
    \item {
        همانطور که در توضیحات نوشته شده، مطلوبیت انتظاری رای‌دهنده
        برای یک رای مانند 
        $B_i$
        برابر است با:
        \[
        \mathbb{E}_{A_{-i}\sim \mu_i(.| A_i)} [u_i(F(B_i, \sigma_{-i}(A_{-i})))]
        \]

        بهترین پاسخ در این مدل بازی 
        TODO
    }
    \item {
        همانطور که در بخش قبلی ثابت کردیم در قاعده
        \lr{AV}
        رای صادقانه همواره بهترین پاسخ است. بنابراین هر بهترین پاسخ
        باید سودی برابر رای صادقانه ایجاد کند. اگر در باور رای‌دهنده
        $i$ 
        حالتی وجود داشته باشد که
        \[
        u_i(F_K^{AV}(B_i, A_{-i})) < u_i(F_K^{AV}(A_i, A_{-i}))
        \]
        و احتمال رخ دادن آن ناصفر باشد، آنگاه $B_i$ 
        نمی‌تواند یک پاسخ بهینه باشد. چرا که مطلوبیت موردانتظار آن اکیدا کمتر از 
        $A_i$
        می‌شود. پس باورهایی که در آنها مطلوبیت 
        $B_i$ 
        و
        $A_{i}$
        برابر است باید در مجموع احتمال یک داشته باشند.

        بنابراین اگر برای هر  رای غیرصادقانه مانند $B_i$
         حالتی ناصفر در باور رای‌دهنده وجود
        داشته باشد که مطلوبیت
        $A_{i}$
        از
        $B_i$
        بیشتر باشد، 
        آنگاه $A_i$ 
        تنها بهترین پاسخ در پروفایل رای صادقانه است.

        حالت‌هایی که در باور رای‌دهنده احتمال $0$
        دارند اهمیتی در تحلیل استراتژی ندارند. اما باقی حالات، هرچند اندک
        ممکن است تاثیرگذار باشند. (به این حالات، حمایت باور او می‌گوییم)
        ‌به طور کلی در بازی‌های بیزی،
         امیدریاضی مطلوبیت هر کنش باید ملاک قرار بگیرد، و ممکن است کنشی
         که در یک حالت نابهینه باشد در مجموع بهترین پاسخ باشد.
         اما در این بازی
        استراتژی وجود دارد که در تمام حالات یکی از پاسخ‌های بهینه است. بنابراین برای اینکه
        یک رای دیگر بتواند از پاسخ‌های بهینه باشد باید در تمام حمایت باور این رای‌دهنده
        مطلوبیتی برابر با استراتژی بهینه داشته باشد.
    }
\end{enumerate}

\section{برگزار‌کننده}
\begin{enumerate}
    \item {
        به ترتیب نشان می‌دهیم چرا هر کدام حالت خاص معیار بعدی خود است.

        \begin{itemize}
            \item {
                فرض کنید کمیته
                $W$
                دارای
                \lr{PJR}
                باشد. آنگاه اگر 
                $l = 1$
                قرار دهیم، نتیجه می‌شود برای هر گروه
                $S \in N$ 
                که
                $S \ge \frac{n}{k}$
                و 
                اشتراک $A_i$ها
                حداقل یک عضو داشته باشد داریم:
                \[
                \left| (\bigcup_{i\in S} A_i) \cap W\right| \ge 1
                \]
                یا معادلا حداقل برای یک 
                $i \in S$
                داشته باشیم:
                \[
                A_i \cap W \neq \emptyset
                \]
                که این دقیقا نشان می‌دهد این کمیته 
                دارای \lr{JR}
                هم هست.

                از طرفی برای $k$
                بزرگتر، می‌توان با تغییر $l$
                نتایج قویتری با \lr{PJR}
                گرفت. 

                مثلا فرض کنید $6$
                رای‌دهنده داشته باشیم، $4$
                کاندید و $k = 3$.
                $4$
                بازیکن به $\{c_1, c_2\}$
                رای دهند، یک بازیکن به $\{c_3\}$
                و یک بازیکن هم به $\{c_4\}$.
                در این بازی، کمیته 
                $\{c_1, c_3, c_4\}$
                دارای 
                \lr{JR}
                است، چرا که
                تمام بازیکنان در کمیته نهایی کاندید موردتایید دارند.
                اما با درنظر گرفتن $4$ 
                نفر اول، تعداد آنها برابر 
                $\frac{2}{3}$
                افراد است، و اشتراک آنها هم برابر $2$ 
                است. اما اشتراک آنها با کمیته نهایی تنها یک است 
                که نشان می‌دهد این کمیته دارای
                \lr{PJR}
                نیست

                این موضوع نشان می‌دهد که \lr{PJR}
                از
                \lr{JR}
                قوی‌تر است.
            }
            \item {
                فرض کنید کمیته $W$
                دارای \lr{EJR}
                باشد. آنگاه هر مجموعه 
                $S$
                از افراد که در شروط 
                \lr{PJR}
                صدق می‌کند چون در شروط 
                \lr{EJR}
                هم صدق می‌کند پس یک عضو مانند 
                $j \in S$ 
                دارند که اشتراک
                $A_j$
                با $W$
                حداقل \lr{l}
                باشد. داریم:
                \[
                \left| (\bigcup_{i\in S} A_i) \cap W\right| \ge
                 \left| A_i \cap W \right| \ge l
                \]
                پس این کمیته در معیار
                \lr{PJR}
                هم صدق می‌کند.

                اما مثال‌هایی هم وجود دارد که در 
                \lr{PJR}
                صدق می‌کنند، ولی در \lr{EJR}
                نه. 
                برای مثال فرض کنید
                $8$
                رای‌دهنده، $7$کاندید و
                $k=4$
                باشد.
                برای هر $1\le i < j \le 4$
                یک بازیکن باشد که مجموعه موردتاییدش برابر
                $\{c_i, c_j, c_5, c_6, c_7\}$
                (در کل $6$ بازیکن)
                و $2$ بازیکن دیگر هم 
                مجموعه
                $\{c_1, c_2, c_3, c_4\}$
                را قبول داشته باشند. در این بازی کمیته
                $\{c_1, c_2, c_3, c_4\}$
                دارای \lr{PJR}
                است. (این موضوع به سادگی قابل بررسی است)
                اما اگر $6$ بازیکن اول را درنظر بگیریم،
                تعداد و اشتراک آنها شرایط
                \lr{EJR}
                برای
                $l = 3$
                را ارضا می‌کند، اما در بین آنها کسی نیست که
                اشتراکش با کمیته نهایی حداقل 
                $3$
                باشد. پس این کمیته دارای
                \lr{EJR}
                نیست.


                بنابراین
                \lr{EJR}
                از
                \lr{PJR}
                قوی‌تر است.
            }
            \item {
                فرض کنید کمیته $W$
                در هسته باشد. فرض خلف می‌کنیم که این کمیته دارای 
                \lr{EJR}
                نباشد. آنگاه
                گروه 
                $S$
                با شرایط گفته شده وجود دارد که در نتیجه
                \lr{EJR}
                صدق نمی‌کند.
                مجموعه $T$
                از کاندیدها را اینگونه در نظر می‌گیریم:
                (با توجه به شروط \lr{EJR} چنین چیزی وجود دارد)

                \[
                |T| = l \quad and \quad T \subseteq \bigcap_{i \in S} A_i 
                \]

                این $T$
                و
                $S$
                در شرط اندازه بلوکه کردن صدق میکنند.
                از آنجایی که $W$
                در هسته است و بلوکه نمی‌شود پس 
                $j \in S$
                وجود دارد که:

                \[
                l = \left| A_j \cap T\right| \le \left|A_j \cap W\right|
                \]

                این موضوع نشان می‌دهد که $A_j$
                اشتراکی بزرگترمساوی از $l$
                دارد
                که خلاف فرض خلف است. پس طبق برهان خلف نتیجه می‌شود که هر 
                مجموعه‌ای که در هسته باشد دارای
                \lr{EJR}
                است.

                اما در هسته بودن خاصیت قوی‌تری است. به طور شهودی دلیل آن این است که 
                گروه‌های بزرگی که لزوما به اندازه بزرگیشان اشتراک ندارند ممکن است
                کمیته را بلوکه کنند، در حالی که اصلا شروط اولیه 
                \lr{EJR}
                را برآورده نمی‌کنند. با یک مثال این موضوع را نشان می‌دهیم:

                        \begin{align*}
            &k=3, n=6, m=4 \\
            &A_1 = \{c_1\} \\ 
            &A_2 = \{c_1, c_2, c_3\} \\ 
            &A_3 = \{c_1, c_2, c_4\} \\ 
            &A_4 = \{c_1, c_2, c_5\} \\ 
            &A_5 = \{c_3, c_4, c_5\} \\ 
            &A_6 = \{c_3, c_4, c_5\} \\
            \\ 
            &W = \{c_3, c_4, c_5\}
        \end{align*}

        در این مثال، کمیته داده شده دارای
        \lr{EJR}
        است. چراکه اگر $S$
        دوعضوی باشد حتما یکی از آنها کاندیدی با شماره بالاتر از یک دارند، و اگر چهار
        عضوی باشد، آنگاه در شرط اولیه 
        \lr{EJR}
        صدق نمی‌کند. اما این کمیته در هسته نیست.
        در نظر می‌گیریم:
        \begin{align*}
            S &= \{1, 2, 3, 4\} \\
            T &= \{c_1, c_2\}
        \end{align*}

        در این صورت به راحتی قابل مشاهده است که $S$
        این کمیته را بلوکه می‌کند.

        درنتیجه در هسته بودن، قوی‌تر از 
        \lr{EJR}
        است.

                
            }
        \end{itemize}
    }
    \item {
        خیر، لزوما برآورده نمی‌کند. به طور شهودی اگر یک گروه از افراد با تعداد
        بیشتر از نصف به یک گروه از کاندیدها رای دهند آنگاه تمام کمیته از همین
        کاندیدها میشود
        و شرط \lr{JR}
        می‌تواند نقض شود.

        به طور دقیق می‌توان این مثال را درنظر گرفت:
        \begin{align*}
            &k=3, n=5, m=4 \\
            &A_1 = \{c_1\} \\ 
            &A_2 = \{c_1\} \\ 
            &A_3 = \{c_2, c_3, c_4\} \\ 
            &A_4 = \{c_2, c_3, c_4\} \\ 
            &A_5 = \{c_2, c_3, c_4\} \\ 
        \end{align*}
        در این مثال کمیته
        $\{c_2, c_3, c_4\}$
        انتخاب می‌شود. اما
        $S = \{1,2\}$
        درحالی که شروط \lr{JR}
        را دارند اما عضوی از آن وجود ندارد که در کمیته نماینده داشته باشد.
        
        بنابراین نتیجه می‌شود قاعده \lr{AV}
        لزوما \lr{JR}
        را برآورده نمی‌کند.
    }
    \item {
        در ابتدا ثابت می‌کنیم قاعده
        \lr{PAV}
        دارای هر سه معیار گفته‌شده می‌باشد. برای اثبات این موضوع طبق
        نکته‌ای که در سوال اول این بخش بررسی کردیم تنها کافی است ثابت کنیم
        که کمیته منتخب در این قاعده دارای
        \lr{EJR}
        است. 
        
        فرض خلف می‌کنیم. فرض کنید کمیته $W$
        انتخاب شده‌باشد که در این قاعده صدق نمی‌کند. در این صورت گروهی
        مانند 
        $S$
        وجود دارد که:

        \[
        |S| \ge \frac{ln}{k}, \qquad \forall i \in S: \: \left| A_i \cap W \ \right| < l
        \]

        در این صورت کاندیدی مانند $c \in \bigcap_{i\in S} A_i$
        وجود دارد که در کمیته نیست. فرض کنید خارج کردن
        کاندید $d$
        از کمیته امتیاز آن کمیته در قاعده 
        \lr{PAV}
        را به اندازه $f(d)$
        کاهش دهد. از آنجایی اگر یک رای‌دهنده به $q$
        کاندید در کمیته رای داده باشد در مجموع
        $f(d)$ها 
        به اندازه $1$
         اثر دارد، چرا که خروج هر کاندید
         به اندازه $\frac{1}{q}$
         از امتیاز موثر او در کمیته کم می‌کند. پس:
         \[
         \sum_{i \in W} f(d) = n
         \]
         بنابراین کاندیدی مانند در $c'$
         در کمیته وجود دارد که:
         \[
         SCORE_{PAV}(W - {c'}) \ge SCORE_{PAV}(W) - \frac{n}{k}
         \]
         در واقع این کاندید، کاندیدی است که $f$
         آن کمینه است. اضافه کردن $c$
         در کمیته باعث می‌شود امتیاز هر کدام از اعضای $S$
         در این کمیته حداقل $\frac{1}{l}$
         افزایش می‌یابد. چراکه هر کدام از آنها قبل از این کمتر از $l$
         عضو در کمیته نهایی داشته‌اند، و اکنون حداقل یک عضو به آن اضافه شده. پس:
         \begin{align*}
         SCORE_{PAV}(W - \{c'\} + \{c\}) &\ge SCORE_{PAV}(W- \{c'\}) + \frac{1}{l} \times \frac{ln}{k} \\
            &\ge SCORE_{PAV}(W) + \frac{n}{k} - \frac{n}{k} \\
            &\ge SCORE_{PAV}(W)
         \end{align*}
 
         درنتیجه کمیته
         $W' = W - \{c'\} + \{c\}$
         امتیاز بیشترمساوی TODO

         برای قاعده
         \lr{ACC}
         ثابت می‌کنیم که دارای
         \lr{JR}
         است، اما دارای
         \lr{PJR}
         نیست.
          
         ابتدا با برهان خلف ثابت می‌کنیم که در معیار \lr{JR}
         صدق می‌کند. فرض کنید 
         $S$ گروهی بیشتر از
         $\frac{n}{k}$
         رای‌دهنده باشند که کاندید $c$
         در اشتراک مجموعه‌های موردتاییدشان باشند اما هیچ‌کدام در کمیته نهایی نماینده
         نداشته باشند. فرض کنید $c'$
         در کمترین مجموعه تاییدی دیگر بازیکنان به عنوان تنها کاندید در کمیته،
         حضور داشته باشد. در این صورت خارج کردن $c'$
         حداکثر امتیاز
         $\frac{n - \frac{n}{k}}{k}$
         را از کمیته کم می‌کند و اضافه کردن $c$
         امتیاز $\frac{n}{k}$ 
         را به کمیته اضافه می‌کند. پس کمیته‌ای که با حذف $c'$
         و اضافه کردن $c$
         به دست می‌آید امتیاز بالاتری دارد که تناقض است. (در این صورت نباید $W$ انتخاب می‌شد.)
         پس کمیته‌ای که با قاعده
         \lr{Approval CC}
         به دست می‌آید دارای 
         \lr{JR}
         است.


    }
    \item {
        هدف برگزارکننده این است که با توجه به رای‌های داده شده
        کمیته‌ای انتخاب شود که هر گروه هم رای، به نسبت اندازه 
        آن گروه در کمیته
         نماینده داشته باشد.
        برگزارکننده لزوما به دنبال گرفتن رای صادقانه از رای‌دهندگان نیست. چرا که
        برای مثال در دنیای واقعی، اگر رای‌دهنده در برگه رای خود صداقت نداشته باشد
        در ادامه هم نمی‌تواند شاکی باشد که چرا مطلوبیت او کم شده. در واقع معیار
        عدالت،
        رای
        رای‌دهنده هاست، نه ترجیح واقعی آنها.

        با توجه به‌این موارد برگزارکننده ترجیح می‌دهد معیارهایی انتخاب کند که از 
        اقلیت‌ها پشتیبانی میکند. در قاعده
        \lr{AV}
        ممکن است، گروهی بزرگ از افراد بدون نماینده بمانند و این ربطی به صداقت یا 
        عدم صداقت آنها ندارد. اما در 
        \lr{Approval CC}
        و
        \lr{PAV}
        تا حدی
        تضمین می‌شود این اقلیت‌‌ها به اندازه کافی نماینده دارند.
    }
    \item {
        یک بار دیگر نمونه را در اینجا می‌آوریم:
        \begin{align*}
            &k=3, n=5, m=4 \\
            &A_1 = \{c_1\} \\ 
            &A_2 = \{c_1\} \\ 
            &A_3 = \{c_2, c_3, c_4\} \\ 
            &A_4 = \{c_2, c_3, c_4\} \\ 
            &A_5 = \{c_2, c_3, c_4\} \\ 
        \end{align*}

        این موضوع که پروفایل رای صادقانه در \lr{AV}
        تعادل نش است را در بخش اول ثابت کردیم. کمیته‌ای که طبق این پروفایل
        و قاعده
        \lr{AV}
        حاصل می‌شود کاندید های $2$
        تا
        $4$
        هستند. ثابت کردیم که این کمیته حتی \lr{JR}
        هم نیست، چه برسد به اینکه در هسته باشد. با اینحال به طور دقیق موارد
        خواسته شده را مشخص می کنیم:
        \begin{align*}
            &S = \{1, 2\} \\
            &T = \{c_1\} \\
            &|T| = 1 \le \frac{|S|}{n}k = \frac{2\times 3}{5} = \frac{6}{5} \\
        \end{align*}
        اما داریم:
        \[
        \forall i \in \{1, 2\},  \quad
    \left| A_i \cap T \right| = 1 > \left| A_i \cap W\right| = 0
        \]
        پس این مجموعه‌ها $W$ را 
        بلوکه می‌کنند.

        این مثال نشان می‌دهد نبود انحراف اجتماعی برای دفاع از اقلیت‌ها کافی نیست. همانطور
        که در این مثال دیدیم هر چند کاندید اول نسبت خوبی از رای‌دهندگان را به عنوان حامی داشت
        اما با توجه به حمایت اکثریت از کاندیدهای دیگر شانسی برای حضور در کمیته $3$
        نفره نداشت. هواداران او هم نمی‌توانستند با استراتژی‌های غیر صادقانه امتیاز
        او را افزایش دهند. دلیل این موضوع این است که قاعده 
        \lr{AV}
        ارزشی برای رای اقلیت قائل نیست، و هر چند موجب صداقت رای‌دهندگان
        می‌شود، اما معیارهای دفاع گروهی را ازدست می‌دهد.


    }
\end{enumerate}


\section{پرسش نهایی}
\subsection*{از کجا شروع کردم؟}
خب، در این قسمت با دو پرسش روبه‌رو شدیم که به همدیگر ربط دارند.
اگر پرسش اول جوابش "خیر" باشد، آنگاه پرسش دوم هم واضحا
"خیر" 
می‌شود. همچنین، بلعکس این موضوع هم درست است، اگر جواب دومین پرسش
"بله"
باشد، آنگاه جواب اولی هم "بله" می‌شود. در نتیجه میشه انتظار داشت
که جواب اولی بله و جواب دومی خیر باشد. (وگرنه
احتمالا منطقی‌تر بود که تنها یکی از آنها در پروژه قرار می‌گرفت!)

از آنجایی که پرسش اولی کمی کلی‌تر است و به نظر می‌رسد کار کردن و تحلیل آن
مشکل‌تر باشد و نیاز به شهود بیشتری داشته‌باشد، در ابتدا به سراغ دومی رفتم.
پرسش دوم ساختار مشخص‌تری دارد و بررسی و تحلیل آن راحت‌تر است.

\subsection*{مثال‌های اولیه برای پرسش دوم}
در ابتدا طبق عادت والبته راهنمایی انجام شده در پروژه چند مثال
کوچک را بررسی کردم. به نظرم رسید می‌شود در ابتدا کم بودن $k$
را بررسی کرد. (در اینجا فکر می‌کردم بررسی 
$n,m$ 
کوچک
چندان نکته‌ای ندارد)

 $k=1$
 یعنی تنها یک کاندید انتخاب می‌شود. با توجه به شرط بلوکه کردن،
 کمیته جایگزین حداکثر یک عضوی است! بدون عضو بودن بدیهی رد شد و
 اگر یک عضو  داشته باشد و در شرط بلوکه کردن صدق کند،
 یعنی یک عضو محبوب برای همه وجود دارد که در کمیته نیست! ضمن اینکه
 کمیته با هیچ رایی اشتراک ندارد! عجب انتخابات باطلی. بدیهی است که
 قاعده
 \lr{PAV}
 همان عضو محبوب را انتخاب می‌کرد. پس برای $k=1$ 
 قاعده همیشه این خاصیت را دارد.

 در $k=2$ حالت‌های 
 $|T| = 0, 2$
 به سادگی رد می‌شوند. برای 
 $|T| =1$ یعنی
 گروهی از رای‌دهندگان به اندازه نصف جمعیت وجود دارد که همه یک کاندید
  را قبول دارند ولی در کمیته نهایی نماینده ندارند. در اینحالت
  ارزش کمیته فعلی حداکثر $\frac{3}{2}n$
  است. در صورتی که اگر به جای کاندید ضعیفتر انتخاب شده
  این کاندید در $T$
  را قرار می‌دادیم امتیاز کمیته حداقل
  $\frac{3}{2}n$
  می‌شد، و اتفاقا کمیته حاصل در هسته قرار می‌گیرد. چرا که دیگر نیمی از جمعیت
  وجود ندارد که بی نماینده باشد.

  در اینجا یک نتیجه جالب هم گرفتم. به‌طور کلی کمیته‌ای که با
  قاعده پاو\footnote{از اینجا به بعد به جهت سادگی \lr{PAV} را همان پاو می‌نویسم}
  انتخاب شود با کمیته جایگزینی به اندازه
  $|T| =1$
  یا 
  $|T| = k$
  بلوکه نمی‌شود. برای اندازه یک، اگر بلوکه شود، یعنی کمیته حتی دارای
  \lr{JR}
  هم نیست. اما طبق چیزهایی که تااینجا بررسی شده، می‌دانیم که 
  کمیته منتخب با قاعده پاو تو این معیار صدق میکنه. پس اندازه
  $T$
  نمی‌تواند یک باشد. اگر $k$
  هم باشد، واضح است که خود $T$
  امتیاز پاو بیشتری از کمیته فعلی دارد، پس برابر $k$
  هم نمی‌تواند باشد.

  با توجه به این توضیحات، در بررسی $k =3$ تنها حالت
  $|T| = 2$ 
  را باید بررسی می‌کردم. این موضوع، یعنی دوسوم بازیکنان هستند که اشتراکشان
  با دو کاندید از کمیته فعلی بیشتر است.
  پس امتیاز پاو کمیته فعلی حداکثر:
  \begin{equation*}
    |S| + SCORE_{PAV} (N / S)
\end{equation*}
  است.
  فرض کنید این $2$ کاندید
  را به جای $2$
   کاندید دیگر در کمیته قرار دهیم.
  دو کاندید حذف شده را میشه طوری انتخاب کرد که
    امتیاز پاو باقی افراد حداکثر به اندازه 
    $n - |S|$
    کاهش یابد. 
    (
        این قسمت مشابه استدلال‌هایی است که در پرسش سوم بخش
        برگزارکننده کردیم.
    )
    پس امتیاز کمیته جدید حداقل برابر است با:

    \[
    \frac{3}{2}|S| + SCORE_{PAV} (N / S)  - (n - |S|)
    \]

    که چون اندازه $S$ حداقل
    دوسوم است این امتیاز از کمیته فعلی بیشتر است.
    (شاید مساوی میشد که با کمی استدلال جزئی‌تر میشد استدلال کرد بیشتر اکید می‌شود)
    پس بازهم گروه بلوکه کننده نداریم و کمیته منتخب این قاعده در
    $k\le 3$
    در هسته است.

    \subsection*{شاید واقعا خروجی پاو در هسته باشد!؟}
    در اینجا گفتم شاید واقعا خروجی این قاعده در هسته باشد و پرسش اول واقعا مقدمه‌ای
    برای این سوال باشد. ایده‌ای که از بررسی $k=3$
    گرفتم، این بود مستقیما $T$
    را در کمیته قرار دهم و یک سری کاندید از آن خارج کنم تا به کمیته‌ای با امتیاز 
    بیشتر برسم و مستقیما حکم را نتیجه بگیرم.
    در اینجا این شهود را داشتم، که اگر این ایده در مثالی با این خواص کار کند
    احتمالا در تمام حالات کار می‌کند:
    (انگار با این فروض به بدترین حالت می‌رویم. هرچند اثبات دقیقی نیست، اما
    برای حرکت به سمت یک اثبات دقیق موثر است.
    اگر هم درست نباشد مارا به سمت یک مثال نقض راهنمایی می‌کند.
    )

    \begin{itemize}
        \item فرض کنیم شرط در حالت تساوی است یعنی $|S| = \frac{n}{k}|T|$
        \item فرض کنیم افزودن $T$ به کمیته باعث می‌شود هر کدام از اعضای $S$
        به ارزششان $\frac{1}{|T|}$
        اضافه شود.
        \item باقی رای‌دهندگان به طور مساوی و تک‌رای به اعضای کمیته داده‌اند
    \end{itemize}

    در این صورت جابجایی $T$
    با بخشی از کمیته فعلی احتمالا امتیاز پاو را به اندازه
    \[
    \frac{n}{k} |T| \times \frac{1}{|T|} = \frac{n}{k}
    \]
    افزایش می‌دهد و به اندازه 
    \[
    \frac{n - \frac{n |T|}{k}}{k} |T|
    \]
    کاهش می‌دهد. (بازیکنان خارج از $S$ که رایشان از کمیته خارج می‌شود.)
    پس برای اینکه نشان دهیم امتیاز پاو این کمیته جدید از قبلی بیشتر است
    باید داشته باشیم:
    \begin{align*}
        \frac{n - \frac{n |T|}{k}}{k} |T| \le \frac{n}{k}
         \Leftrightarrow (1 - \frac{|T|}{k}) |T| \le 1
    \end{align*}

    اما این عبارت در $|T| = \frac{k}{2}$
    بیشترین حالت خودش می‌شود. (با مشتقگیری به این نتیجه رسیدم)
    که در آن بسیار از $1$
    بیشتر می‌شود.

    هر چند این ایده جواب نداد، اما الان فهمیده‌ایم که برای ساخت
    یک مثال نقض، تقریبا چه ساختاری را باید رعایت کنیم.

    \subsection*{ساختار یک مثال خوب}
    احتمالا در مثالی که خروجی قاعده پاو این خاصیت را نداشته باشد
    باید 
    نیمی از بازیکنان به
    $\frac{k}{2}$
    کاندید انتخاب نشده 
    علاقه داشته باشند و باقی رای‌دهندگان 
    یا به همه کمیته فعلی، یا به صورت یکنواخت به یک عضو کمیته فعلی علاقه داشته
    باشند.
    نوشتن شروط نقض برای $k$ 
    کلی پیچیده و طولانی شد و با توجه به قاعده پاو که عبارات
    $H_i$
    در آن ظاهر می‌شود به نظر رسید محاسبات قرار نیست سرانجامی داشته باشد.

    \subsection*{کمی تفکر روی پرسش اول}
    در اینجا کمی ناامید شدم و مدتی فکر روی مسئله را کنار گذاشتم.
    در بازگشت، به این نتیجه رسیدم شاید خوب باشد کمی پرسش اول را بررسی کنم
    چرا که به هرحال این احتمال وجود داشت که جواب آن خیر باشد و پرسش دوم بیهوده!
    همچنین با بررسی‌های انجام شده در پرسش دوم، این حس در من تقویت شد که ممکن است مثال‌هایی
    وجود داشته باشد که کمیته‌ای در هسته نباشد. در واقع، دو مسیر برای اثبات پرسش اول
    (در صورت وجود آن)
    به ذهنم می‌رسید:
    \begin{itemize}
        \item  طبق قاعده‌ای کمیته‌ای انتخاب شود که حتما در هسته باشد.
        \item یکجوری ثابت کنیم بین تمام حالات انتخاب کمیته، حداقل یکی از آنها در هسته است.
    \end{itemize}

    سعی کردم با تغییرات جزئی در قاعده پاو، آن را بهبود ببخشم تا کمیته منتخب
    آن حتما در هسته باشد، اما به جایی نرسیدم. در واقع صرفا اندکی وزن امتیازدهی این 
    قاعده را تغییر دادم که به نظر خوب نرسید. مثلا به جای 
    $\frac{1}{i}$
    برای $i$امین کاندید یک بازیکن
    وزن را
    $\frac{1}{\sqrt{i} }$
    یا
    $\frac{c}{i}$
    قرار دادم. اما چیز واضحی به چشم نیامد.

    روش دوم هم بسیار کلی به نظر می‌رسید و کار کردن روی آن سخت. 

    \subsection*{ بررسی چند مثال بیشتر از ساختار خوب }





\end{document}